\chapter{Introduction}
\label{chap:introduction}

\section{Context and Motivation}

Large language models are large. A model with seven billion parameters in float16 requires roughly fourteen gigabytes of storage and a similar amount of RAM or VRAM to hold the weights alone. In serverless and multi-tenant inference environments, GPU memory is a shared resource. When a model is not actively serving requests, keeping it resident wastes memory that could serve other tenants. Systems therefore evict idle models and reload them on demand.

When a request arrives for an evicted model, the inference pipeline cannot begin until the checkpoint is fully loaded into memory. For large models this loading step takes time. In a production serving environment, this manifests as cold-start latency: the delay between a request arriving and the model being ready to generate the first token. For latency-sensitive applications this delay is directly visible to users.

Modern storage devices offer substantial raw bandwidth. Despite this, checkpoint loading in practice often achieves a fraction of the available throughput. The gap between peak hardware capability and effective loading performance suggests that the bottleneck lies not only in the storage medium itself, but also in how loading requests are submitted, batched, and integrated into the surrounding software stack.

This project investigates the checkpoint loading path as a systems problem. It examines how backend choice affects throughput and latency for real LLM checkpoints, how that effect changes with checkpoint layout and model size, and how much of the backend-level gain survives once Python bindings and a serving engine are included.

\section{Checkpoint Loading as an I/O Submission Problem}

Checkpoint loading is fundamentally a read-intensive workload. A loader issues a set of read operations derived from checkpoint metadata. Each operation specifies a file path, a byte offset, and a length. Tensor boundaries and offsets must be respected, which constrains how freely a loader can reorder or merge requests.

Storage hardware, particularly NVMe devices with internal queueing, performs best when given many outstanding requests to work with concurrently. The internal parallelism of the device can only be exploited if the submission path sustains enough useful in-flight work. Synchronous I/O, where each thread blocks waiting for its current read before issuing the next, can limit queue depth. Asynchronous submission can sustain more outstanding work, but only if the backend architecture is designed well enough to keep the machine busy. This turns out to be an important caveat in the final results.

The central question of this project is therefore not simply whether asynchronous I/O is better than synchronous I/O. The more useful question is whether LLM checkpoint loading can be driven by adaptive heuristics: using a simpler synchronous path when the model is small enough that extra submission machinery is not worthwhile, and switching to io\_uring once model size and shard structure make deeper parallelism pay off.

\section{TensorStore}

This project presents TensorStore, a checkpoint loading framework designed to evaluate I/O submission strategies under controlled and reproducible conditions. TensorStore parses checkpoint metadata and produces a normalised set of byte-range read requests that are then executed by a selectable I/O backend. By keeping the request trace identical across backends, the framework isolates the effect of submission strategy from checkpoint layout or parsing overhead.

TensorStore supports two checkpoint formats. SafeTensors stores all tensor data contiguously in a single file with a metadata header. This layout produces a largely sequential access pattern once the header is parsed. The ServerlessLLM-style partitioned layout distributes tensor ranges across multiple files, increasing request count and file-switching overhead. This layout places greater stress on the submission path and better exposes the differences between I/O models.

Three eager-loading backends are implemented. The synchronous backend issues blocking POSIX read calls. The async backend uses Tokio-based grouped loading. The io\_uring backend uses low-level Linux io\_uring with batched submission and worker-owned rings. The project also exposes mmap-backed lazy-access APIs for Python and vLLM experiments where they are useful.

Python bindings are provided to enable integration testing within representative serving workflows. Integration with vLLM is used as a concrete case study.

\section{Contributions}

This project makes three concrete contributions. The first is a reproducible evaluation framework, TensorStore, that isolates backend effects from checkpoint parsing and layout by generating comparable loading workloads across backends. The second is a systematic experimental comparison of synchronous, async, and io\_uring loading across real LLM checkpoints and two checkpoint formats, quantifying where each backend is strongest and where it is not. The third is an analysis of how backend choice interacts with checkpoint layout, model size, Python binding overhead, and vLLM integration, leading to a practical adaptive heuristic for checkpoint loading --- in broad terms, sync for smaller models and io\_uring for larger ones, with format-specific nuances around that rule.

\section{Report Structure}

The report proceeds as follows. Chapter~\ref{chap:background} reviews related work in checkpoint formats, serverless inference systems, and Linux I/O mechanisms. Chapter~\ref{chap:system_context} defines the system context, describes the cold-start pipeline, and clarifies the experimental constraints. Chapter~\ref{chap:design} describes the TensorStore architecture and design decisions. Chapter~\ref{chap:implementation} covers the implementation and software engineering process. Chapter~\ref{chap:methodology} describes the experimental methodology. Chapter~\ref{chap:results} presents the results. Chapter~\ref{chap:discussion} discusses the findings. Chapter~\ref{chap:professional_issues} examines professional and ethical considerations. Chapter~\ref{chap:conclusion} concludes.
